% ----------------------------------------------------------------
% Section: System Architecture
% ----------------------------------------------------------------
\chapter{SYSTEM ARCHITECTURE}
\label{ch:architecture}

\section{Overview}

The \texttt{meta-rpi5-secure-aosp} project delivers a complete, reproducible pipeline for building and deploying a security-hardened Android Automotive OS image onto a Raspberry Pi~5. The system has two primary concerns: a \textit{build-time pipeline} that produces a signed, encrypted SD card image, and a \textit{runtime boot chain} that verifies and mounts partitions before handing off to Android userspace.

\section{Boot Chain}
\label{sec:arch_bootchain}

The secure boot chain consists of five layers, each responsible for initializing and handing off to the next:

\begin{enumerate}
    \item \textbf{RPi Firmware} (VideoCore, closed-source): Powers on, initializes DRAM and PCIe, reads \texttt{config.txt} from the FAT32 boot partition, and loads \texttt{u-boot.bin}. The firmware also provides the Device Tree blob to U-Boot via the ARM64 boot protocol \cite{rpi_hardware_docs}.

    \item \textbf{U-Boot}: Runs the compiled boot script (\texttt{boot\_avb.scr}). Invokes \texttt{avb init} and \texttt{avb verify} to check the VBMeta structure against the compiled-in root public key. On success, loads the Android kernel \texttt{Image} and \texttt{ramdisk.img} from the FAT32 partition and calls \texttt{booti} \cite{uboot_build_gcc, uboot_rpi5_forum}.

    \item \textbf{Android Kernel} (Linux 6.12, \texttt{android\_rpi5\_defconfig}): Decompresses and initializes hardware. Diagnostic markers (\texttt{RPi5-DIAG:}) are emitted via \texttt{pr\_notice()} at four boundaries: \texttt{start\_kernel} entry, after \texttt{console\_init}, at \texttt{do\_initcalls}, and before \texttt{run\_init\_process}.

    \item \textbf{Android Init (first-stage)}: Mounts \texttt{/system} and \texttt{/vendor} using the \texttt{fstab} embedded in the ramdisk. The \texttt{avb=vbmeta} flag triggers dm-verity verification of both partitions via the \texttt{vbmeta} partition \cite{android_dm_verity}.

    \item \textbf{Android Init (second-stage)} and userspace: Loads SELinux policy, starts \texttt{vold} (which activates FBE on \texttt{/data}), and launches the Android Automotive framework \cite{android_fbe}.
\end{enumerate}

Figure~\ref{fig:boot_chain} illustrates the complete boot flow across the five coloured layers described above. Each layer either advances control along the green \textit{Yes} path or routes to a device reboot along the red \textit{No} path. The three diamond-shaped decision nodes represent the cryptographic verification gates: the AVB signature check performed by U-Boot against its compiled-in root public key; the dm-verity block-integrity check applied to \texttt{/system} and \texttt{/vendor} at kernel mount time; and the FBE key-derivation check performed by \texttt{vold} before \texttt{/data} becomes accessible. All three failure paths converge on a single \textit{Boot Failure / Reboot} terminal, enforcing a fail-closed policy with no degraded-mode fallback.

\begin{figure}[H]
\centering
\begin{tikzpicture}[
    node distance=0.5cm,
    startstop/.style={
        rounded rectangle,
        draw=black!70, fill=green!25,
        minimum width=3.4cm, minimum height=0.65cm,
        align=center, font=\small\bfseries
    },
    process/.style={
        rectangle,
        draw=black!70, fill=blue!10,
        minimum width=3.4cm, minimum height=0.65cm,
        align=center, font=\small
    },
    decision/.style={
        diamond,
        draw=black!70, fill=orange!20, aspect=2.5,
        align=center, font=\small
    },
    failnode/.style={
        rounded rectangle,
        draw=red!80!black, fill=red!15,
        minimum width=3.2cm, minimum height=0.65cm,
        align=center, font=\small\bfseries
    },
    arrow/.style={->, >=Stealth, thick},
    failarrow/.style={->, >=Stealth, thick, draw=red!70!black}
]

%% Main vertical flow
\node[startstop]                     (A)  {Start};
\node[process,   below=0.40cm of A]  (B)  {ROM Boot};
\node[process,   below=0.25cm of B]  (C)  {Secondary Bootloader};

\node[process,   below=0.80cm of C]  (D)  {U-Boot Loaded};
\node[process,   below=0.25cm of D]  (E)  {AVB Validation};
\node[decision,  below=0.45cm of E]  (F)  {AVB Valid?};

\node[process,   below=0.55cm of F]  (G)  {Kernel Loaded};
\node[process,   below=0.25cm of G]  (H)  {dm-verity};
\node[decision,  below=0.45cm of H]  (I)  {dm-verity Pass?};

\node[process,   below=0.55cm of I]  (J)  {Early AOSP Init};
\node[process,   below=0.25cm of J]  (K)  {FBE Validation};
\node[decision,  below=0.45cm of K]  (L)  {/data Mounted?};
\node[startstop, below=0.55cm of L]  (M)  {AOSP Running};

%% Failure node (right of kernel layer)
\node[failnode, right=4.2cm of I]    (FAIL) {Boot Failure\\/ Reboot};

%% Layer background boxes
\begin{scope}[on background layer]
    \node[draw=gray!50,   fill=gray!8,   rounded corners=4pt, inner sep=5pt,
          fit={(A)(B)(C)}]        (bg_hw)   {};
    \node[draw=blue!50,   fill=blue!5,   rounded corners=4pt, inner sep=5pt,
          fit={(D)(E)(F)}]        (bg_ub)   {};
    \node[draw=cyan!60,   fill=cyan!5,   rounded corners=4pt, inner sep=5pt,
          fit={(G)(H)(I)}]        (bg_kn)   {};
    \node[draw=purple!40, fill=purple!4, rounded corners=4pt, inner sep=5pt,
          fit={(J)(K)(L)(M)}]     (bg_ao)   {};
    \node[draw=red!40,    fill=red!5,    rounded corners=4pt, inner sep=5pt,
          fit={(FAIL)}]           (bg_fail) {};
\end{scope}

%% Layer labels (right of each background box at the north edge)
\node[font=\scriptsize\itshape, anchor=west, text=gray!60!black]
    at ($(bg_hw.north east)+(0.12,0)$)   {Hardware / Firmware};
\node[font=\scriptsize\itshape, anchor=west, text=blue!70!black]
    at ($(bg_ub.north east)+(0.12,0)$)   {Bootloader (U-Boot)};
\node[font=\scriptsize\itshape, anchor=west, text=cyan!70!black]
    at ($(bg_kn.north east)+(0.12,0)$)   {Kernel};
\node[font=\scriptsize\itshape, anchor=west, text=purple!60!black]
    at ($(bg_ao.north east)+(0.12,0)$)   {Android Userspace (AOSP)};
\node[font=\scriptsize\itshape, anchor=south, text=red!70!black]
    at ($(bg_fail.north)+(0,0.10)$)      {Failure Handling};

%% Main flow arrows
\draw[arrow] (A) -- (B);
\draw[arrow] (B) -- (C);
\draw[arrow] (C) -- (D);
\draw[arrow] (D) -- (E);
\draw[arrow] (E) -- (F);
\draw[arrow] (F) -- node[right, font=\scriptsize] {Yes} (G);
\draw[arrow] (G) -- (H);
\draw[arrow] (H) -- (I);
\draw[arrow] (I) -- node[right, font=\scriptsize] {Yes} (J);
\draw[arrow] (J) -- (K);
\draw[arrow] (K) -- (L);
\draw[arrow] (L) -- node[right, font=\scriptsize] {Yes} (M);

%% No-branch arrows to FAIL
%% F (above FAIL): go right, then route down to FAIL.north
\draw[failarrow] (F.east)
    -- node[above, font=\scriptsize] {No}
    ++(0.55,0) coordinate (fR)
    -- (fR |- FAIL.north) -- (FAIL.north);
%% I (same level): direct right to FAIL
\draw[failarrow] (I.east)
    -- node[above, font=\scriptsize] {No}
    (FAIL.west);
%% L (below FAIL): go right, then route up to FAIL.south
\draw[failarrow] (L.east)
    -- node[above, font=\scriptsize] {No}
    ++(0.55,0) coordinate (lR)
    -- (lR |- FAIL.south) -- (FAIL.south);

\end{tikzpicture}
\caption{Secure boot chain on Raspberry Pi~5. Three verification gates (AVB, dm-verity, FBE) each enforce a fail-closed policy: any verification failure triggers a device reboot with no fallback to a degraded boot state.}
\label{fig:boot_chain}
\end{figure}

\section{Partition Layout}
\label{sec:arch_partitions}

The GPT-partitioned SD card is defined in \texttt{config/sdcard/uboot\_aosp\_signed.yaml}. The Table~\ref{tab:partitions} describes the layout.

\begin{table}[h]
\centering
\caption{GPT Partition Layout (signed build, 14.6~GB SD card)}
\label{tab:partitions}
\begin{tabular}{cllrl}
\toprule
\textbf{No.} & \textbf{Name} & \textbf{Format} & \textbf{Size (MB)} & \textbf{Contents} \\
\midrule
p1 & \texttt{boot}     & FAT32 & 128    & RPi firmware, U-Boot, kernel, ramdisk, \texttt{boot.scr} \\
p3 & \texttt{userdata} & ext4  & (fill) & \texttt{/data} --- FBE-encrypted user data \\
p5 & \texttt{system}   & ext4  & 3072   & \texttt{/system} --- AVB hashtree signed \\
p6 & \texttt{vendor}   & ext4  & 384    & \texttt{/vendor} --- AVB hashtree signed \\
p7 & \texttt{metadata} & ext4  & 16     & \texttt{/metadata} --- FBE key material \\
p8 & \texttt{vbmeta}   & raw   & 4      & AVB top-level VBMeta image \\
\bottomrule
\end{tabular}
\end{table}

The non-contiguous partition numbers (p1, p3, p5--p8) are intentional: p2 and p4 are reserved for alternative configurations without breaking existing \texttt{fstab} or U-Boot environment references.

\section{Build Pipeline}
\label{sec:arch_pipeline}

The build pipeline is orchestrated by \texttt{rpi5-build}, a Python CLI implemented in \texttt{src/meta\_rpi5\_secure\_aosp/}. It is driven by a YAML configuration file (\texttt{config/rpi5\_uboot\_aosp\_signed.yaml}) and accepts per-run overrides via CLI flags.

\subsection{Stages}

\begin{description}
    \item[\texttt{sync}] Clones U-Boot at a pinned commit, initializes the AOSP repo with the \texttt{android-16.0.0\_r4} manifest, and applies local manifest overlays from \texttt{ganeshhalthota/android\_local\_manifest} \cite{android_local_manifest}. Optionally clones the kernel source from \texttt{raspberry-vanilla/android\_kernel\_brcm\_rpi} at a pinned commit.

    \item[\texttt{patch}] Applies \texttt{git format-patch} stacks from \texttt{patches/uboot/}, \texttt{patches/kernel/}, and \texttt{patches/aosp/device\_brcm\_rpi5/} using \texttt{git am --3way}. Idempotent: already-applied patches are detected via \texttt{git log --grep} on the commit subject, preventing double-application.

    \item[\texttt{build}] Compiles U-Boot (\texttt{rpi\_arm64\_defconfig}), enables AVB Kconfig symbols (\texttt{CONFIG\_LIBAVB}, \texttt{CONFIG\_AVB\_VERIFY}, \texttt{CONFIG\_CMD\_AVB}), injects the AVB root public key into \texttt{avb\_verify.c}, compiles the boot script (\texttt{boot\_avb.scr}), and builds AOSP (\texttt{bootimage}, \texttt{systemimage}, \texttt{vendorimage}).

    \item[\texttt{sign}] Signs \texttt{system.img} and \texttt{vendor.img} with AVB hashtree footers (SHA-256/RSA-4096), shrinks images if needed to fit within the predefined partition sizes, and generates the combined \texttt{vbmeta.img} \cite{avb2_readme}.

    \item[\texttt{sdcard}] Assembles the signed images and supporting files into a flashable GPT image using the partition layout YAML.
\end{description}

\subsection{Configuration Model}

The build is governed by a two-layer configuration model:

\begin{itemize}
    \item \textbf{rpi5 config YAML} (\texttt{config/rpi5\_uboot\_aosp\_signed.yaml}): Controls source tags, AVB key paths, signing behaviour, SELinux mode, encryption mode, cmdline profile, and boot state.
    \item \textbf{sdcard config YAML} (\texttt{config/sdcard/uboot\_aosp\_signed.yaml}): Defines the GPT partition layout, sizes, image paths, and extra file mappings.
\end{itemize}

A \texttt{BuildContext} dataclass carries the resolved configuration across all stage modules, ensuring a single source of truth for all options. Security-sensitive options (\texttt{boot\_state\_override = orange}, \texttt{avb\_fail\_policy = fail\_open}) are blocked by default and require the explicit \texttt{--allow-insecure-boot-state} flag.

\section{Security Design Decisions}
\label{sec:arch_security}

\subsection{External Signing Pipeline}

The AVB signing is deliberately separated from the AOSP build
(\texttt{BOARD\_AVB\_ENABLE = false}). This design:
\begin{itemize}
    \item Keeps the private key outside the AOSP build environment.
    \item Allows the same AOSP images to be signed with different keys (development vs. production) without rebuilding.
    \item Matches the typical automotive production workflow where signing is a separate,
    audited step.
\end{itemize}

\subsection{Fail-Closed AVB Policy}

The default \texttt{avb\_fail\_policy = fail\_closed} causes U-Boot to call \texttt{reset} if AVB initialization or verification fails. This prevents a compromised image from booting in a degraded state and is the recommended policy for production devices \cite{avb2_readme}.

\subsection{FBE Dependency Chain}

FBE requires a strict prerequisite chain enforced by the build tool at startup:
\begin{enumerate}
    \item \texttt{sdcard.enable\_signing = true} (AVB must be active).
    \item \texttt{avb.uboot\_fail\_policy = fail\_closed}.
    \item \texttt{userdata} and \texttt{metadata} partitions present in the sdcard config.
\end{enumerate}
If any prerequisite is missing, the build tool raises a descriptive error before any work is performed, preventing a partially-configured image from being produced.
